% !TEX root = ../main.tex
% =====================================================================
% 3. METHOD -- page budget 2.00 (including the architecture figure)
%
% Order: the baseline first (so the reader knows what we are changing),
% then the encoder, then the bottleneck, then the head, then the cost
% consequence. Only the equations the contribution needs; the other six
% encoder variants and the other five pooling heads are in Appendix B.
% =====================================================================

\section{Method}
\label{sec:method}

\begin{figure}[t]
  \centering
  \resizebox{\textwidth}{!}{% Architecture figure, drawn in TikZ so it stays sharp and anonymous.
% Two panels:
%   (a) the whole pipeline: input -> encoder -> per-step embeddings -> head
%       -> prediction AND interpretation (both from the same quantity).
%   (b) the bottleneck: what we changed inside the encoder block.
% Kept deliberately plain -- only the parts the contribution needs.

\begin{tikzpicture}[
    font=\scriptsize,
    node distance=3.5mm,
    box/.style   ={draw, rounded corners=1.2pt, minimum height=7.5mm,
                   align=center, inner sep=2pt, fill=black!3},
    enc/.style   ={box, fill=blue!8},
    head/.style  ={box, fill=orange!12},
    outbox/.style={box, fill=green!10, minimum height=6.5mm},
    tag/.style   ={inner sep=1pt, align=center},
    ar/.style    ={-{Latex[length=1.4mm]}, thin},
  ]

  % ---------------- panel (a): the pipeline -------------------------
  \node[box, minimum width=11mm] (in) {$\mathbf{x}$\\[-1pt]$1{\times}T$};
  \node[enc, right=of in, minimum width=11mm] (stem) {stem\\[-1pt]conv};
  \node[enc, right=of stem, minimum width=24mm] (blocks)
        {$L\times$ multi-scale separable\\[-1pt]block, dilation $1{\dots}16$};
  \node[box, right=of blocks, minimum width=13mm] (emb)
        {$\mathbf{H}$\\[-1pt]$d{\times}T$};

  \node[head, right=6mm of emb, yshift=5.5mm, minimum width=17mm] (attn)
        {attention $a_j^k$};
  \node[head, right=6mm of emb, yshift=-5.5mm, minimum width=17mm] (clf)
        {per-step logits $\hat{y}_j^{\,k}$};

  \node[circle, draw, inner sep=0.6pt, right=4mm of emb, xshift=21mm] (mul) {$\times$};
  \node[box, right=3mm of mul, minimum width=15mm, fill=red!8] (ev)
        {evidence $g_j^k$\\[-1pt]$C{\times}T$};

  \node[outbox, right=5mm of ev, yshift=5.5mm, minimum width=25mm] (pred)
        {\textbf{prediction}: mean of the $k$ largest $g_j^k$};
  \node[outbox, right=5mm of ev, yshift=-5.5mm, minimum width=25mm] (interp)
        {\textbf{interpretation}: $\mathbf{I}_{k,j} = g_j^k$};

  \draw[ar] (in)     -- (stem);
  \draw[ar] (stem)   -- (blocks);
  \draw[ar] (blocks) -- (emb);
  \draw[ar] (emb.east) -- ++(2mm,0) |- (attn.west);
  \draw[ar] (emb.east) -- ++(2mm,0) |- (clf.west);
  \draw[ar] (attn.east) -| (mul.north);
  \draw[ar] (clf.east)  -| (mul.south);
  \draw[ar] (mul) -- (ev);
  \draw[ar] (ev.east) -- ++(1.5mm,0) |- (pred.west);
  \draw[ar] (ev.east) -- ++(1.5mm,0) |- (interp.west);

  \begin{scope}[on background layer]
    \node[draw, dashed, rounded corners, inner sep=3pt, fill=blue!3,
          fit=(stem)(blocks)] (encbox) {};
    \node[draw, dashed, rounded corners, inner sep=3pt, fill=orange!4,
          fit=(attn)(clf)(mul)(ev)] (headbox) {};
  \end{scope}
  \node[tag, above=0.5mm of encbox] {\textbf{encoder} (length preserving)};
  \node[tag, above=0.5mm of headbox] {\textbf{Top-$k$ conjunctive pooling head}};

  % ---------------- panel (b): the bottleneck -----------------------
  \coordinate (bstart) at ($(in.west) + (0,-19mm)$);

  \node[anchor=west, tag] (blabel) at (bstart)
        {\textbf{inside one block:}\\[1pt] the pointwise mix};

  \node[box, right=5mm of blabel, minimum width=9mm] (b1) {$d$};
  \node[box, right=6mm of b1, minimum width=9mm] (b2) {$d$};
  \draw[ar] (b1) -- node[above, tag] {$1{\times}1$} node[below, tag] {$d^{2}$} (b2);

  \node[anchor=west, tag, right=9mm of b2] (barrow) {becomes};

  \node[box, right=5mm of barrow, minimum width=9mm] (c1) {$d$};
  \node[box, right=5.5mm of c1, minimum width=6mm, fill=red!10] (c2) {$r$};
  \node[box, right=5.5mm of c2, minimum width=9mm] (c3) {$d$};
  \draw[ar] (c1) -- node[above, tag] {squeeze} (c2);
  \draw[ar] (c2) -- node[above, tag] {expand} (c3);
  \node[tag, below=1mm of c2] {$2dr = d^{2}/2$ at $r = d/4$};

\end{tikzpicture}
}
  \caption{\emph{Top:} the model. The encoder never changes the length, so every downstream quantity
  keeps a time axis. The head's two branches multiply into the evidence $g_j^k$, the single quantity
  from which \emph{both} outputs are read: the prediction is a reduction of it over time, the
  interpretation \emph{is} it. \emph{Bottom:} the change inside the block --- the pointwise $1\times1$
  channel mix is factorised through a narrow middle layer, halving the term that dominates the
  parameter count.}
  \label{fig:arch}
\end{figure}

\subsection{The baseline we start from}
\label{sec:method:mil}

A series $\bag \in \R^{1 \times \nsteps}$ is treated as a bag of $\nsteps$ instances, one per time
step. An encoder $f$ produces $\mathbf{H} = f(\bag) \in \R^{\dmodel \times \nsteps}$, one embedding
$\emb_j$ per step. \millet{}'s \emph{conjunctive pooling} then runs two branches over those
embeddings and multiplies them:
\begin{equation}
  a_j = \frac{\exp\!\big(w^{\!\top}\tanh(V \emb_j)\big)}{\sum_{i} \exp\!\big(w^{\!\top}\tanh(V \emb_i)\big)},
  \qquad
  \instpred_j = W_{\mathrm{c}} \emb_j \in \R^{\nclz},
  \qquad
  g_j^k = a_j \, \instpred_j^{\,k},
  \label{eq:conj_gate}
\end{equation}
where $a_j$ says how much step $j$ matters and $\instpred_j^{\,k}$ says what step $j$ argues for
regarding class $k$. The product $g_j^k$ is the \emph{evidence} of step $j$ for class $k$, and both
outputs are read from it:
\begin{equation}
  \bagpred^{\,k} = \frac{1}{\nsteps}\sum_{j=1}^{\nsteps} g_j^k
  \qquad\text{(prediction)},
  \qquad\qquad
  \interp_{k,j} = g_j^k
  \qquad\text{(interpretation)} .
  \label{eq:conj_bag}
\end{equation}
This is why the explanation is faithful by construction: it is not a model of the model, it is the
summand of the prediction. Three properties of \Cref{eq:conj_gate,eq:conj_bag} matter later.
\textbf{(L1)} the gate $a_j$ is shared across classes, so every class is forced to look in the same
place; \textbf{(L2)} the normalisation in $a_j$ is over time, so attention mass is conserved rather
than absolute; and \textbf{(L3)} the reduction is a fixed mean over all $\nsteps$ steps, so a signal
confined to a short window is divided by the length of the whole series.

\paragraph{Where the cost is.}
The MIL machinery in \Cref{eq:conj_gate} is almost free. Replacing \millet{}'s interpretable head with
ordinary global average pooling on the same InceptionTime backbone changes the model from 423{,}707
to 422{,}666 parameters: explaining itself costs $0.25\,\%$ of the model, and everything else is the
encoder. A cheaper interpretable classifier is therefore a cheaper \emph{encoder} problem, not a
pooling problem --- which is what \Cref{sec:method:enc} addresses, while \Cref{sec:method:head}
addresses quality at a fixed cost.

\subsection{A length-preserving separable encoder}
\label{sec:method:enc}

We replace InceptionTime with a stack of \emph{multi-scale separable residual blocks}. Each unit
reads its input with three depthwise convolutions of kernel size $\{5, 11, 23\}$ at a shared
dilation, sums them, and mixes channels with a pointwise $1\times1$ convolution, followed by
BatchNorm \citep{ioffe2015batchnorm}, ReLU and dropout; two units and a residual connection form a
block, and dilation grows geometrically across blocks up to a cap of 16. Two design constraints are
forced by the interpretation rather than by accuracy: convolutions use ``same'' padding with no
stride and no pooling, so $\nsteps$ is preserved at every stage --- required, because
\Cref{eq:conj_bag} needs one value per input step and the faithfulness metrics perturb individual
steps --- and the dilation cap keeps every embedding tied to a bounded window, so importance stays
local and readable. The third choice is the cost one: the depthwise/pointwise split
\citep{chollet2017xception, howard2017mobilenets} replaces the $\dmodel^2 k$ weights of a dense
convolution with
\begin{equation}
  \underbrace{\dmodel k}_{\text{depthwise}} \;+\; \underbrace{\dmodel^{2}}_{\text{pointwise}}
  \;\ll\; \dmodel^{2} k ,
  \label{eq:sep_cost}
\end{equation}
which is what makes the small setting reachable at all: at $\dmodel = 128$ a dense convolution
covering the same three kernel sizes would cost $638{,}976$ weights per step against $22{,}144$ for
the separable unit.

\paragraph{The bottleneck.}
\Cref{eq:sep_cost} has a consequence: once the convolutions are separable, the \emph{pointwise} term
dominates. At $\dmodel = 128$ the three depthwise convolutions of one unit cost $5{,}376$ weights
while the single $1\times1$ costs $16{,}512$ --- three times all three depthwise convolutions
combined. We therefore factorise that $1\times1$ through a narrow middle layer of width
$r = \dmodel / 4$, squeezing and then expanding:
\begin{equation}
  P = P_{\mathrm{expand}}\, P_{\mathrm{squeeze}}, \qquad
  P_{\mathrm{squeeze}} \in \R^{r \times \dmodel}, \quad
  P_{\mathrm{expand}} \in \R^{\dmodel \times r},
  \qquad \text{cost } \dmodel^{2} \rightarrow 2\dmodel r = \dmodel^{2}/2 .
  \label{eq:bottleneck}
\end{equation}
The construction is the standard squeeze-and-expand bottleneck \citep{he2016resnet,
iandola2016squeezenet, sandler2018mobilenetv2}; what matters here is where it is applied and what it
leaves intact. It halves the dominant term independently of $\dmodel$, and the block still consumes
and returns $(\dmodel, \nsteps)$, so length preservation --- and with it the per-step interpretation
--- is untouched. At the setting used throughout ($\dmodel = 64$, 4 blocks) it takes the complete
model from 57{,}580 to 41{,}324 parameters and from $109.7$ to $76.7$\,MFLOPs per series: $28\,\%$
of the parameters and $30\,\%$ of the arithmetic, for a change that is invisible to the head.

\subsection{Top-$k$ conjunctive pooling}
\label{sec:method:head}

The head is where limitations L1 and L3 are addressed, at zero parameter cost. First we give each
class its own gate, $a_j \rightarrow a_j^k$, so that classes may disagree about where to look (L1).
Then, instead of the fixed mean of \Cref{eq:conj_bag}, we average only the largest evidence values
for each class, with the count set as a fraction $\kappa$ of the series length so that the head is
length-agnostic:
\begin{equation}
  k = \max\big(1, \lceil \kappa \nsteps \rceil\big),
  \qquad
  \bagpred^{\,k} = \frac{1}{k}\!\!\sum_{j \in \mathcal{S}^k_{\mathrm{top}}}\!\! g_j^k,
  \qquad
  \interp_{k,j} = g_j^k ,
  \label{eq:topk}
\end{equation}
where $\mathcal{S}^k_{\mathrm{top}}$ indexes the $k$ largest $g_j^k$. Selection of the strongest
instances is not new in MIL \citep{durand2016weldon}; the form here is per-class, expressed as a
fraction of length, and applied to conjunctive evidence, which gives it a property the sweep in
\Cref{sec:res:ablation} depends on.

\begin{property}[Exact reduction]
\label{prop:topk}
At $\kappa = 1$, \Cref{eq:topk} is class-wise conjunctive pooling term for term. For $\kappa < 1$ the
gradient reaches only the selected steps.
\end{property}

Two consequences. Practically, $\kappa$ is a hyperparameter and not a weight, so this head has
\emph{exactly} the same parameter count as the head it generalises --- the interpretation quality it
buys is free in every cost column of \Cref{tab:cost}. Experimentally, the $\kappa = 1$ row of
\Cref{tab:ablation} is not a different model but the same model with selection switched off, so any
difference along that column is attributable to selection and to nothing else.

\paragraph{Training objective.}
Every model is trained with cross-entropy on the bag logits plus one term that exists for the
explanation rather than for the classification:
\begin{equation}
  \mathcal{L} = \mathrm{CE}(\bagpred, y)
  \;+\; \lambda_{\mathrm{focus}} \Big({-}\textstyle\sum_{j} \bar a_j \log(\bar a_j + \epsilon)\Big),
  \label{eq:loss}
\end{equation}
with $\bar a_j$ the class-averaged attention. Attention spread evenly over a series has high entropy,
so adding entropy to a minimised loss pushes the model towards concentrated attention. Without it a
model can be right while its importance map is flat, and a reader cannot tell where the evidence
stops. We set $\lambda_{\mathrm{focus}} = 0.01$ for our models and $0$ for the \millet{} baseline,
since the term is our addition; \Cref{sec:res:ablation} reports what it is worth.
